\section{SCENARIO \#11: 'Devastation of the Lower Valley": August 16-23,
}


\begin{scenariodata}
\textbf{Game Length:} 8 Turns\\
\end{scenariodata}

\begin{scenariodata}
(Weather: July - October Column)\\
\end{scenariodata}

\begin{scenariodata}
\textbf{First Phase:} Union\\
\end{scenariodata}

\begin{scenariodata}
\textbf{Union:}\\
\end{scenariodata}

\begin{scenariodata}
Resource Factor 211\\
\end{scenariodata}

\begin{scenariodata}
\textbf{Initial Placement:}\\
\end{scenariodata}

\begin{scenariodata}
Hex A-16 (Romney) : (41, 1A], (1F, 1S)\\
\end{scenariodata}

\begin{scenariodata}
B\&O : (151, 4A, 2C\}, (4F, 3S)\\
\end{scenariodata}

\begin{scenariodata}
Hex S-28 (Front Royal) -- :\\
\end{scenariodata}

\begin{scenariodata}
\textbf{and adjacent :} 6C, 2H, (15, 1W)\\
\end{scenariodata}

\begin{scenariodata}
Hex N-20 (Winchester) 2 SI, (1F, 3W)\\
\end{scenariodata}

\begin{scenariodata}
Hexes N-24, M-23, : 611, 10A, 3C, 1H,\\
\end{scenariodata}

\begin{scenariodata}
\textbf{and adjacent :} (Sheridan, Wright,\\
\end{scenariodata}

Emory, Crook, 2S, 2W)

\begin{scenariodata}
Hex Z-13 : Available Turn \#2\\
\end{scenariodata}

\begin{scenariodata}
Adds + 1 Resource Factor : 5C, 2H\\
\end{scenariodata}

\begin{scenariodata}
\textbf{Withdrawal Hexes:}\\
N-7, S-10, Z-13, CC-15 through CC-33\\
\end{scenariodata}

\begin{scenariodata}
(Supply/Wagon Entry Hexes:\\
N-7, \$-10, Z-13, A-11, A-13, A-16, CC-15 through CC-33\\
\end{scenariodata}

\begin{scenariodata}
One additional Fortification counter available.\\
\end{scenariodata}

\begin{scenariodata}
\textbf{Confederates:}\\
\end{scenariodata}

\begin{scenariodata}
\textbf{Resource Factor :} 14\\
\end{scenariodata}

\begin{scenariodata}
\textbf{Initial Placement:}\\
\end{scenariodata}

\begin{scenariodata}
Hexes L-26, M-27, N-28 : 441,10A, 7C, 2H,\\
\end{scenariodata}

\begin{scenariodata}
\textbf{and adjacent :} (Early, Rodes, Gordon,\\
Breckinridge, Ramseur,\\
4F, 58,5W)\\
\end{scenariodata}

\begin{scenariodata}
Hex CC-40 : 111, 4A, 7C, 2H,\\
\end{scenariodata}

\begin{scenariodata}
\textbf{and adjacent :} (Kershaw, 18, 1W)\\
\end{scenariodata}

\begin{scenariodata}
Hex 0-51 (Staunton) 214i, (1F, 1S, 1W)\\
\end{scenariodata}

\begin{scenariodata}
Hex BB-53 (Charlottesville); 21, 1A, (1F, 1S)\\
\end{scenariodata}

\begin{scenariodata}
(Partisans: Mosby, Gilmor, McNeill)\\
\end{scenariodata}

\begin{scenariodata}
\textbf{Withdrawal Hexes:}\\
\end{scenariodata}

\begin{scenariodata}
CC-40 through CC-58\\
\end{scenariodata}

\begin{scenariodata}
(Supply/Wagon Entry Hexes:\\
\end{scenariodata}

\begin{scenariodata}
CC-40 through CC-58, O-51, M-50, E-46)\\
\end{scenariodata}

\begin{scenariodata}
No additional Fortification counters.\\
\end{scenariodata}

\begin{scenariodata}
\textbf{Victory Conditions:} All other results are a draw.\\
\end{scenariodata}

\begin{scenariodata}
\textbf{ADVANCED GAME:} Union wins if they have 15 or more lead in total\\
\end{scenariodata}

\subsection*{Victory Points.}

\begin{scenariodata}
Confederates win if they have 10 or more lead in\\
\end{scenariodata}

total Victory Points.

\begin{scenariodata}
\textbf{BASIC GAME:} Same as the Advanced Game, only "spot'' the Con-\\
\end{scenariodata}

federates 5 Victory Points at the start of the game.

\begin{scenariodata}
\textbf{HISTORICAL RESULTS:} Draw.\\
\end{scenariodata}

\begin{scenariodata}
\textbf{PLAYTESTER'S COMMENTS:} The Union player may attempt a full\\
\end{scenariodata}

scale attack on Early's forces, but the general weakness in supplies and the unfavorable attrition rate against Fortifications does not encourage this course of action. Also, Kershaw's units bring the Confederates up to a rough parity in strength with the Union forces. The Union's best bet is to start devastation, and try to defeat some segment of the Confederate army (especially Kershaw's division). The Confederates can afford to fight more than in most scenarios, but equal attrition could cost the game. Watch for any careless Union maneuvers that could open up Harper's Ferry for a quick dash.

a

On hearing of the approach of the Confederate reinforcements, and still awaiting reinforcements of his own, Sheridan decided to fall back again towards Harper's Ferry. During this withdrawal, which the

\begin{scenariodata}
Confederates found themselves unable to do much about, the Union\\
\end{scenariodata}

forces systematically devastated the areas they moved through, thus denying these resources to the Confederates. The Rebels were again able to move north and cut the B\&O Railroad.

Both commanders came under criticism for their actions in these series of maneuvers. Sheridan was criticized for allowing the Confederates to again cut the B\&O. Early was censured for allowing the Union forces to devastate the Lower Valley. This increased both sides' desire for a final show-down battle in the Shenandoah.

Sheridan continued to receive reinforcements, and to re-organize until the middle of September, while the press continued to howl for action. Finally, U.S. Grant paid a personal visit to the Valley, and gave Sheridan the orders to 'Go In!" Sheridan was already prepared to move, and, hearing that Kershaw's division had left the Valley, and that the over-confident Early had spread his army over a wide area, the Union army moved into attack positions...

